\section*{Introduction}
\phantomsection\addcontentsline{toc}{section}{Introduction}

The previous chapter showed that global structural features carry useful information for
DR grading, and that graph and topological signals should be fused rather than simply
concatenated. This chapter turns these observations into a complete grading framework,
TOPORET~\cite{belhadj2025plos}. TOPORET connects patients in an inductive population graph
whose edges combine two similarities: how alike the fundus images look to a CNN, and how
alike their vessel networks are topologically. This thesis introduces it as a DR grading
system that addresses the relational limitation (L1) and the structural limitation (L2)
at the same time; we are not aware of an earlier system combining these two elements in a
single inductive population graph.

We first review related work on population graphs and topological features in medical
imaging. We then state the two hypotheses tested in this chapter, present the six-stage
pipeline and the dual-similarity graph, and report the experiments: statistical tests of
the hypotheses, analysis of the graph, ablation, cross-dataset evaluation and comparison
with the literature. We close with a discussion of what the results do and do not imply.

\section{Related Work}
\label{sec:ch3-related}

\paragraph{Population graphs in medical imaging.}
Representing a cohort as a graph whose nodes are patients was popularised for disease
prediction from brain imaging, where graph convolutional networks combine imaging features
with phenotypic similarity between subjects~\cite{parisot2018disease}. The idea is that
similar patients tend to share a diagnosis, so the graph can propagate information from
labelled to unlabelled subjects. These approaches are usually transductive: the whole
cohort, including the test patients, must be present when the model is trained, which is
not compatible with a screening service where new patients arrive every day. In the DR
grading work reviewed here, graph-based approaches have used visual similarity only; we re-implement such a
visual-only graph (GraphDR) as the control of this chapter.

\paragraph{Topological features in medical imaging.}
Persistent homology has been used to describe the shape of anatomical structures and to
feed learning models with stable, interpretable descriptors~\cite{hofer2017topological,
carriere2020perslay}. For the retina, the complexity of the vascular tree, often summarised
by its fractal dimension~\cite{masters2004fractal}, is known to change with disease. To our
knowledge, persistence descriptors of the vessel skeleton had not been used as edge
information in a population graph before this work.

\section{Problem Formulation and Hypotheses}
\label{sec:ch3-problem}\label{sec:hyp}

Let $\{(I_j,y_j)\}$ be a set of fundus images with ICDR grades $y_j\in\{0,\ldots,4\}$. The
task is to predict the grade $\hat y_i$ of a new image $I_i$ that was not seen during
training, using both its own content and the content of similar patients. The chapter tests
two pre-specified hypotheses; their decision thresholds were fixed before any data were
analysed.

\paragraph{H1 -- The Relational Hypothesis.}
An inductive population graph connecting patients through a dual-similarity metric
(visual embedding similarity $+$ Wasserstein topological distance) provides a
statistically significant QWK improvement over independent per-image classification.
\emph{Pre-specified test:} paired $t$-test across five cross-validation folds, TOPORET
versus EfficientNet-B3; null hypothesis $\mu_{\Delta\mathrm{QWK}} \leq 0$. The nominal
one-tailed critical value is $t_{0.05,4}=2.132$; the more conservative $t_4 > 2.78$
($t_{0.025,4}$) was adopted as the decision threshold.

\paragraph{H2 -- The Structural Hypothesis.}
$H_0$ and $H_1$ persistent-homology features of the retinal vascular skeleton are
associated with ICDR severity grade.
\emph{Pre-specified test:} one-way ANOVA across the five grades ($n=5{,}000$), Tukey HSD
post-hoc, Bonferroni-corrected significance level $\alpha=0.005$; critical value
$F_{0.005;\,4,4995}=3.72$.

\section{Proposed Approach}
\label{sec:ch3-approach}

\subsection{Overview}
\label{sec:ch3-overview}

TOPORET~\cite{belhadj2025plos} addresses L1 with a population graph and L2 with
persistent homology, in a single pipeline.
\Cref{fig:toporet} shows its six stages, the typed interface between consecutive stages,
and its headline results.

Three design principles guide TOPORET. First, the graph is \emph{inductive}: a new patient
is connected to the existing cohort at inference time, without retraining, which is what a
screening service requires. Second, topology is used twice: as a feature of each patient
and inside the similarity that decides which patients are neighbours, which extends the
synergy observed in \cref{ch:found} from a single image to the whole population. Third, each
stage has a typed input and output, so that it can be tested, replaced or audited on its
own; this modularity is reused unchanged by DOTS in \cref{ch:dots}.

% Figure 3.1 --- TOPORET six-stage pipeline
\begin{figure}[!htbp]
\centering
\resizebox{\linewidth}{!}{%
\begin{tikzpicture}[
  font=\footnotesize,
  box/.style={draw, rounded corners=2pt, align=center, inner sep=2.5pt,
              minimum height=1.6cm, text width=1.58cm},
  arr/.style={-{Latex[length=2.2mm]}, thick}
]
\node[box, fill=black!6, draw=black!55] (s1) {\textbf{S1}\\CLAHE $+$\\skeleton};
\node[box, fill=black!6, draw=black!55, right=2.2mm of s1] (s2) {\textbf{S2}\\EffNet-B3\\(frozen)};
\node[box, fill=axisSf, draw=axisS, right=2.2mm of s2] (s3) {\textbf{S3}\\Rips\\TDA 6-D};
\node[box, fill=fusionf, draw=fusion, right=2.2mm of s3] (s4) {\textbf{S4}\\Fusion\\1542-D};
\node[box, fill=axisRf, draw=axisR, right=2.2mm of s4] (s5) {\textbf{S5}\\Dual-sim.\\graph ($\tau$)};
\node[box, fill=axisRf, draw=axisR, right=2.2mm of s5] (s6) {\textbf{S6}\\{\scriptsize GraphSAGE}\\$+$ softmax};
\foreach \a/\b in {s1/s2,s2/s3,s3/s4,s4/s5,s5/s6} \draw[arr] (\a) -- (\b);
\node[left=2.5mm of s1, font=\footnotesize, align=center] (in) {Fundus\\image $I_i$};
\node[right=2.5mm of s6, font=\footnotesize, align=center] (out) {Grade\\$\hat y_i$};
\draw[arr] (in) -- (s1); \draw[arr] (s6) -- (out);
\node[below=6mm of s3.south east, anchor=north, font=\footnotesize, align=center] {%
  Kaggle DR: QWK $=0.88$, accuracy $95.5\%$ (five-fold cross-validation)\\H1: $t_4=3.82$ \quad H2: $F=84.3$, $p<0.001$};
\end{tikzpicture}}
\caption[TOPORET six-stage pipeline]{TOPORET six-stage pipeline. Blue stages form the relational axis (L1), the green
stage the structural axis (L2). There is no ordinal-safety axis: every image is graded
and the order of the grades is not modelled. Stage contracts are given in
\cref{tab:stages}; results from~\cite{belhadj2025plos}.}
\label{fig:toporet}
\end{figure}


\subsection{Six-Stage Pipeline}
\label{sec:stages}

\Cref{tab:stages} gives the typed input/output contract of each stage. The contracts are
what makes each stage unit-testable in isolation (\cref{sec:iec62304}).

\begin{table}[H]
\centering
\caption[TOPORET six-stage modular pipeline]{TOPORET six-stage modular pipeline. Each stage has a typed interface contract
that supports IEC~62304 unit verification.}
\label{tab:stages}
\small
\begin{tabular}{@{}c l P{4.6cm} P{4.4cm}@{}}
\toprule
Stage & Name & Input $\to$ Output & Key operation \\
\midrule
S1 & Pre-processing & Fundus image $\to$ enhanced image $+$ skeleton & CLAHE $+$ Zhang--Suen thinning \\
S2 & CNN extraction & Enhanced image $\to$ $x_i^{\mathrm{CNN}}\in\mathbb{R}^{1536}$ & EfficientNet-B3 (frozen) \\
S3 & TDA persistence & Skeleton points $\to$ $x_i^{\mathrm{TDA}}\in\mathbb{R}^{7}$ & GUDHI Vietoris--Rips, \cref{eq:tda} \\
S4 & Feature fusion & $x_i^{\mathrm{CNN}}, x_i^{\mathrm{TDA}} \to z_i\in\mathbb{R}^{1543}$ & Concatenation \\
S5 & Graph construction & $\{z_i\}\to G=(V,E,S_{ij})$ & FAISS $k$-NN, dual similarity \\
S6 & GNN inference & $G, z_i \to \hat y_i$ & GraphSAGE $+$ CORN ordinal head \\
\bottomrule
\end{tabular}
\end{table}

\Cref{fig:skeleton} follows stages S1--S3 on one fundus image. Panel~(a) is the segmented
vessel network; panel~(b) its one-pixel skeleton $P_i$, obtained by Zhang--Suen
thinning~\cite{zhang1984thinning}; panel~(c) the persistence diagram $D_i$ of the
Vietoris--Rips filtration built on the skeleton points, in which filled dots are $H_0$
components and open dots are $H_1$ loops. The 7-D vector of \cref{eq:tda} is computed
from this diagram.

\begin{figure}[!tb]
\centering
\includegraphics[width=0.86\textwidth]{vessel_skeleton_persistence}
\caption[Vessel skeleton extraction and topological features (stages S1--S3)]{Vessel skeleton extraction and topological features (stages S1--S3):
(a)~vessel network, (b)~Zhang--Suen skeleton $P_i$, (c)~Vietoris--Rips persistence
diagram $D_i$ (filled: $H_0$ components; open: $H_1$ loops).}
\label{fig:skeleton}
\end{figure}

\Cref{fig:clahe} shows the first operation of stage S1 on a Grade~2 image: the photograph
as acquired and the same image after CLAHE~\cite{zuiderveld1994clahe} with clip limit
$2.0$ and an $8\times 8$ tile grid. Local contrast enhancement sharpens the boundaries of
microaneurysms and haemorrhages and separates the vessels from the background. This makes
the vessel segmentation and skeleton extracted later in stage S1 more reliable, without
saturating healthy tissue or shifting the global colour balance.

\begin{figure}[H]
\centering
\begin{minipage}{0.86\textwidth}
\begin{subfigure}[t]{0.4\textwidth}\centering
  \includegraphics[width=\linewidth,height=\linewidth,keepaspectratio]{clahe_original}
  \caption{Original, as acquired}
\end{subfigure}\hfill
\begin{subfigure}[t]{0.4\textwidth}\centering
  \includegraphics[width=\linewidth,height=\linewidth,keepaspectratio]{clahe_enhanced}
  \caption{After CLAHE (clip 2.0, $8\times8$)}
\end{subfigure}
\end{minipage}
\caption[Effect of CLAHE pre-processing on a Grade~2 fundus image]{Effect of CLAHE pre-processing (stage S1) on a Grade~2 fundus image: (a)~original, (b)~after CLAHE with clip limit $2.0$ and an $8\times8$ tile grid.}
\label{fig:clahe}
\end{figure}

\subsection{Dual-Similarity Population Graph}
\label{sec:dualsim}

The dual-similarity edge weight combines the cosine similarity of visual embeddings with
a Wasserstein topological term:
\begin{equation}
S_{ij} = \alpha\cdot\cos\bigl(x_i^{\mathrm{CNN}},\, x_j^{\mathrm{CNN}}\bigr)
       + (1-\alpha)\cdot\exp\bigl(-W_2(D_i, D_j)\bigr),
\label{eq:dualsim}
\end{equation}
with $\alpha=0.65$ and $k=15$ neighbours, both selected by five-fold cross-validation.
\Cref{tab:hparam} reports the QWK stability of each hyperparameter over its range: every
hyperparameter moves QWK by at most $0.5$~pp inside the stable interval, so the results
are not a knife-edge consequence of one setting.

\begin{table}[H]
\centering
\caption[Hyperparameter sensitivity of TOPORET]{Hyperparameter sensitivity of TOPORET: QWK stability over the tested range.}
\label{tab:hparam}
\small
\begin{tabular}{@{}lllc@{}}
\toprule
Hyperparameter & Range tested & QWK stability & Selected \\
\midrule
$\alpha$ (visual/topological balance) & $[0.3,\,0.9]$ & $\pm0.5$~pp over $[0.5,\,0.8]$ & 0.65 \\
$k$ (neighbours) & $\{5,8,10,15,20\}$ & $\pm0.3$~pp over $[8,\,15]$ & 15 \\
$d_h$ (hidden dimension) & $\{128,256,512\}$ & $\pm0.2$~pp & 256 \\
$p$ (dropout) & $[0.1,\,0.5]$ & $\pm0.4$~pp over $[0.2,\,0.5]$ & 0.30 \\
\bottomrule
\end{tabular}
\end{table}

\Cref{alg:toporet} summarises how TOPORET builds the population graph and grades a new
patient. Training images are processed once; a new image only requires its own features,
a $k$-nearest-neighbour query and one GraphSAGE forward pass.

\begin{algorithm}[!htb]
\caption{TOPORET: population-graph construction and inductive grading.}
\label{alg:toporet}
\DontPrintSemicolon
\KwIn{training images $\{I_j\}$ with grades $\{y_j\}$; new image $I_i$; parameters $\alpha=0.65$, $k=15$}
\KwOut{predicted grade $\hat y_i$}
\BlankLine
\ForEach{training image $I_j$}{
  $\tilde I_j \leftarrow \mathrm{CLAHE}(I_j)$; $P_j \leftarrow \mathrm{Skeleton}(\tilde I_j)$\tcp*{stage S1}
  $x_j^{\mathrm{CNN}} \leftarrow \mathrm{EfficientNetB3}(\tilde I_j)$\tcp*{stage S2}
  $D_j \leftarrow \mathrm{RipsPersistence}(P_j)$; $x_j^{\mathrm{TDA}} \leftarrow \phi(D_j)$\tcp*{stage S3, \cref{eq:tda}}
  $z_j \leftarrow [\,x_j^{\mathrm{CNN}} \Vert x_j^{\mathrm{TDA}}\,]$\tcp*{stage S4}
}
Build the graph $G$: connect each node to its $k$ most similar nodes under $S_{ij}$ (\cref{eq:dualsim})\tcp*{stage S5}
Train GraphSAGE with a CORN ordinal head on $G$ (five-fold cross-validation)\tcp*{stage S6}
\BlankLine
Compute $z_i$ for the new image $I_i$ with stages S1--S4\;
$\mathcal{N}(i) \leftarrow$ the $k$ training nodes with highest $S_{ij}$ (FAISS query)\;
$h_i \leftarrow \mathrm{GraphSAGE}(z_i,\{z_j\}_{j\in\mathcal{N}(i)})$\tcp*{\cref{eq:sage-agg,eq:sage-upd}}
\Return{$\hat y_i \leftarrow \mathrm{CORN}(h_i)$}
\end{algorithm}

\subsection{Design Choices}
\label{sec:ch3-design}

Several alternatives were possible at each stage of TOPORET. This subsection explains the
choices made and the reasons behind them.

\paragraph{Inductive rather than transductive learning.}
Population-graph methods in medical imaging are often transductive: the test patients are
part of the graph during training~\cite{parisot2018disease}. This is convenient for a
research study but incompatible with a screening service, where new patients arrive every
day and retraining the model for each of them is not possible. GraphSAGE was chosen
because it learns an aggregation \emph{function} rather than node-specific embeddings, so
that a new patient can be graded by connecting it to the existing graph and applying the
learned function (\cref{sec:sage}).

\paragraph{Topology computed on vessel branch points.}
Persistent homology can be computed directly on the image intensities, using cubical
complexes, or on a point cloud extracted from the image. The first option captures every
bright or dark structure, including lesions, illumination gradients and noise. Computing
the Vietoris--Rips filtration on the branch points of the vessel skeleton restricts the
analysis to the object of interest, the vascular network, and produces features that can
be interpreted in vascular terms (\cref{sec:tda-bio}). The price is a dependence on the
quality of the segmentation and thinning, which is mitigated by CLAHE pre-processing and
by the stability of persistence diagrams (\cref{eq:stability}).

\paragraph{A small vector of summary statistics.}
Richer representations of persistence diagrams exist, such as persistence landscapes or
images, and learned layers can process diagrams directly (\cref{sec:ph}). A seven-dimensional
vector of summary statistics was preferred for two reasons: each component has a direct
clinical interpretation, which matters for transparency, and the small dimension keeps the
topological signal from being drowned by the $1536$ dimensions of the CNN embedding when
the two are concatenated.

\paragraph{Similarity fused at graph-construction time.}
The hybrid model of \cref{ch:found} concatenated graph and topological signals and found
their gains synergistic. TOPORET goes further and uses topology inside the edge weights
(\cref{eq:dualsim}), so that two patients are connected only if their images look alike
\emph{and} their vessel networks have a similar shape. The Wasserstein distance was chosen
to compare diagrams because it takes all points of the diagrams into account, not only the
largest difference as the bottleneck distance does. The balance $\alpha$ between the two
terms was selected by cross-validation rather than fixed a priori.

\paragraph{Approximate nearest neighbours.}
Connecting each patient to its $k$ most similar patients requires a nearest-neighbour
search among tens of thousands of embeddings. An exact search would grow linearly with the
cohort; the FAISS library~\cite{johnson2021faiss} provides approximate search with
negligible loss of accuracy and a query time of about a millisecond, which keeps the
per-patient cost of TOPORET independent of the cohort size. Together, these choices keep TOPORET simple enough to be verified stage by stage, which prepares the engineering work of \cref{ch:reg}.

\section{Experimentation}
\label{sec:ch3-exp}

\subsection{Datasets and Protocol}

All experiments use Kaggle EyePACS~\cite{kaggle2015eyepacs} ($N=88{,}702$) for training
and in-distribution evaluation, with five folds stratified by ICDR grade and a fixed random
seed. Out-of-distribution generalisation is measured on Messidor-2~\cite{decenciere2014messidor}
and IDRiD~\cite{porwal2018idrid}, which come from other countries and cameras. The ANOVA of
H2 uses a stratified sample of $5{,}000$ images ($1{,}000$ per grade). Three systems are
compared throughout: EfficientNet-B3 alone, the visual-only graph (GraphDR) and TOPORET.
The hyperparameters were selected on the training folds only (\cref{tab:hparam}).

\subsection{Results: Topological Features and Severity (H2)}
\label{sec:h2}

\paragraph{Pre-specified test.}
One-way ANOVA on each of the six TDA features of \cref{tab:anova} across the five ICDR
grades ($n=5{,}000$, 1{,}000 images per grade, stratified from Kaggle EyePACS), with
Tukey HSD post-hoc comparisons~\cite{tukey1949} and a Bonferroni-corrected significance
level. The Bonferroni level for six features is $0.05/6\approx0.0083$; the more
conservative $\alpha=0.005$ was pre-specified, with critical value
$F_{0.005;\,4,4995}=3.72$.

\begin{table}[H]
\centering
\caption[ANOVA of the topological features across ICDR grades]{One-way ANOVA of the topological features across the five ICDR grades
($n=5{,}000$, Kaggle EyePACS). Values are mean $\pm$ SD.}
\label{tab:anova}
\small
\setlength{\tabcolsep}{4pt}
\begin{tabular}{@{}lccccccc@{}}
\toprule
Feature & G0 & G1 & G2 & G3 & G4 & $F(4,4995)$ & $\eta^2$ \\
\midrule
$\Pi_{H_0}$ & $24.3\pm3.1$ & $23.1\pm3.4$ & $21.6\pm3.7$ & $19.8\pm4.0$ & $18.7\pm4.2$ & $385.8^{***}$ & 0.24 \\
$\Pi_{H_1}$ & $8.2\pm1.4$ & $9.4\pm1.8$ & $11.2\pm2.3$ & $13.5\pm3.1$ & $15.6\pm3.8$ & $1308.9^{***}$ & 0.51 \\
$n_{H_0}$ & $142\pm18$ & $136\pm19$ & $127\pm21$ & $118\pm23$ & $109\pm25$ & $388.8^{***}$ & 0.24 \\
$\beta_1$ & $1.8\pm0.6$ & $2.3\pm0.8$ & $3.1\pm1.0$ & $4.0\pm1.3$ & $4.8\pm1.5$ & $1258.4^{***}$ & 0.50 \\
$D_f$ & $1.71\pm0.04$ & $1.68\pm0.05$ & $1.64\pm0.05$ & $1.60\pm0.06$ & $1.57\pm0.06$ & $1177.5^{***}$ & 0.49 \\
$\Delta$ & $0.8\pm0.3$ & $1.1\pm0.4$ & $1.5\pm0.6$ & $2.2\pm0.9$ & $3.8\pm1.3$ & $2294.2^{***}$ & 0.65 \\
\bottomrule
\end{tabular}
\tabnote{$^{***}p<0.001$. Critical value $F_{0.005;\,4,4995}=3.72$: every feature exceeds it
by a factor of at least $100$.}
\end{table}

All six features exceed the critical value, with large to very large effect sizes
($\eta^2\in[0.24,\,0.65]$; \cref{tab:anova}). Maximum $H_1$ persistence
($\Delta$, $\eta^2=0.65$) shows the strongest association, followed by the pre-specified
primary feature, total $H_1$ persistence ($\Pi_{H_1}$, $\eta^2=0.51$), the first Betti number
($\beta_1$, $\eta^2=0.50$) and the fractal dimension ($D_f$, $\eta^2=0.49$), consistent with neovascularisation (Grade~4) as the
hallmark of proliferative DR. In practical terms, an $\eta^2$ of $0.51$ means that the ICDR grade alone explains about half of the variance of the total $H_1$ persistence across images.

\paragraph{H2 supported.}
$F(4,4995)=1308.9$, $p<0.001$, $\eta^2=0.51$ (large effect) for the pre-specified primary
feature $\Pi_{H_1}$. The null hypothesis is rejected with a margin of about $\times 352$
over the Bonferroni-corrected critical value. The analysis demonstrates a strong
statistical association between the extracted topological descriptors and ICDR severity;
it does not show that topology causes, or mechanistically explains, disease severity.

\subsection{Results: Contribution of the Population Graph (H1)}
\label{sec:h1}

\paragraph{Pre-specified test.}
Paired $t$-test across the five cross-validation folds, TOPORET versus EfficientNet-B3
(independent per-image classification). Decision threshold $t_4>2.78$ (conservative;
\cref{sec:hyp}). \Cref{tab:h1} reports the result for the three systems.

\begin{table}[H]
\centering
\caption[H1 confirmatory test: graph versus independent classification]{H1 confirmatory test: QWK gain of a population graph over independent
classification (Kaggle EyePACS, five folds).}
\label{tab:h1}
\small
\begin{tabular}{@{}lcccc@{}}
\toprule
System & QWK (5-fold) & $t_4$ & $p$ (one-tailed) & Cohen's $d_z$ \\
\midrule
EfficientNet-B3 (no graph) & $0.840\pm0.012$ & -- & -- & -- \\
GraphDR (visual graph) & $0.850\pm0.014$ & 1.71 & 0.08 (n.s.) & 0.76 \\
TOPORET (dual sim.\ $+$ TDA) & $0.880\pm0.009$ & 3.82 & 0.009 & 1.71 \\
\bottomrule
\end{tabular}
\end{table}

\paragraph{H1 supported.}
$t_4=3.82$, $p=0.009$, $d_z=1.71$ (very large effect). The population graph is associated
with a $+4.0$~pp QWK gain beyond independent classification, with a margin of $\times1.37$
over the conservative critical value. Because the inferential unit is the cross-validation
fold, the paired $t$-test has four degrees of freedom. Bootstrap and exact permutation
analyses were therefore reported as complementary robustness checks rather than as
independent replications.

\paragraph{Statistical robustness of H1.}
Five folds are a small sample for a paired $t$-test, even though $N=88{,}702$ images are
used for training and evaluation. Three supplementary analyses address this concern:
(i)~a bootstrap CI on $\Delta$QWK ($10{,}000$ resamples of the fold differences) gives
$\Delta\mathrm{QWK}=+4.0$~pp, 95\% CI $[+2.1,\,+5.9]$~pp, which excludes zero;
(ii)~the gain is positive in every fold, with no fold showing a near-zero or negative
difference; (iii)~an exact sign-flip permutation test over all $2^5=32$ sign assignments of
the fold differences therefore gives $p_{\mathrm{perm}}=1/32\approx0.031$, the smallest value
attainable with five folds, which agrees with the $t$-test.

\paragraph{Where the gain comes from.}
\Cref{fig:toporetcm} compares the normalised confusion matrices of the CNN baseline and of
the topology-aware population graph on Kaggle DR~\cite{belhadj2025plos}. The diagonal is
stronger for every grade, and the largest improvement is at the Mild--Moderate boundary,
the most ambiguous one for human graders: Mild images predicted as Moderate drop from $6\%$
to $4\%$, and Moderate images predicted as Mild from $6\%$ to $5\%$. No~DR reaches a
per-class accuracy of $0.97$ and proliferative DR $0.95$. This is the behaviour expected
from a population graph: a borderline image is connected to similar, confidently graded
patients, whose labels help to resolve the ambiguity.

\begin{figure}[!tb]
\centering
\includegraphics[width=\linewidth]{toporet_confusion}
\caption[Confusion matrices of the CNN baseline and the population graph]{Normalised
confusion matrices on Kaggle DR: CNN baseline (left) and topology-aware population graph
(right). Correct predictions lie on the diagonal; the largest reduction of errors is at the
Mild--Moderate boundary. Adapted from~\cite{belhadj2025plos}.}
\label{fig:toporetcm}
\end{figure}

\subsection{Population Graph Structure Analysis}
\label{sec:graphstruct}

TOPORET reaches QWK $=0.880$ on Kaggle EyePACS (in distribution) and $0.900$ on
Messidor-2 (out of distribution), a $+2.0$~pp cross-dataset gain consistent with the
stability of persistence diagrams (\cref{eq:stability}). However, TOPORET fails the
predefined safety threshold: Sev-NR $=3.4\%>2\%$. This residual gap is addressed in
\cref{ch:dots}.

\Cref{tab:graphstruct} compares the visual-only graph with the dual-similarity graph. The
dual metric raises the fraction of same-grade edges from $63.2\%$ to $71.4\%$, and the
Grade~3 neighbourhood purity -- at the highest-stakes boundary -- from $61.3\%$ to $68.7\%$.
Every point of Grade~3 purity means fewer Severe-NPDR patients whose nearest neighbours
carry a different, lower grade, which directly reduces the probability of a
misclassification through the relational signal.

\begin{table}[H]
\centering
\caption[Population-graph structure: visual-only versus dual similarity]{Population-graph structure: visual-only similarity (GraphDR) versus dual
similarity (TOPORET).}
\label{tab:graphstruct}
\small
\begin{tabular}{@{}lcc@{}}
\toprule
Property & Visual only (GraphDR) & Dual sim.\ (TOPORET) \\
\midrule
Same-grade edges & 63.2\% & 71.4\% ($+8.2$~pp) \\
Adjacent-grade edges ($|\Delta|=1$) & 29.8\% & 23.1\% \\
Cross-boundary edges ($|\Delta|\geq2$) & 7.0\% & 5.5\% \\
Grade-3 neighbourhood purity & 61.3\% & 68.7\% ($+7.4$~pp) \\
Grade-4 neighbourhood purity & 79.1\% & 82.3\% ($+3.2$~pp) \\
\bottomrule
\end{tabular}
\end{table}

The gain in purity is largest at Grade~3, where visual similarity alone most often links
severe cases to moderate ones; adding topology to the edge weights moves these patients
closer to neighbours with the same grade.

\subsection{Ablation Study}
\label{sec:ablation3}

\Cref{tab:ablation3} isolates the contribution of each component of TOPORET. The combined
gain ($+0.040$) exceeds the sum of the individual contributions ($+0.016+0.010=+0.026$)
by a factor of $1.54$: an interaction effect consistent with synergy between the dual
graph and the TDA features. The Wasserstein component improves neighbourhood purity,
which in turn amplifies the discriminative value of TDA through message passing.

\begin{table}[H]
\centering
\caption[Sequential ablation of TOPORET]{Sequential ablation of TOPORET on Kaggle EyePACS (five-fold mean).}
\label{tab:ablation3}
\small
\begin{tabular}{@{}lccc@{}}
\toprule
Configuration & QWK & $\Delta$ & Sev-NR \\
\midrule
EfficientNet-B3 alone (baseline) & 0.840 & -- & 4.3\% \\
$+$ TDA features only (no graph) & 0.856 & $+0.016$ & 4.1\% \\
$+$ Visual-only graph (no TDA) & 0.850 & $+0.010$ & 4.1\% \\
$+$ Dual graph $+$ TDA (TOPORET) & 0.880 & $+0.040$ & 3.4\% \\
\bottomrule
\end{tabular}
\end{table}

This experiment differs from the Part~B ablation of \cref{ch:found}
(\cref{tab:hybrid}), which used an intra-image graph and concatenated the topological
descriptors as a separate channel. Here the graph links patients and the dual metric uses
both sources jointly to define its edges, and the synergy between them is again observed.

\subsection{Cross-Dataset Generalisation}
\label{sec:crossdata}

\Cref{tab:crossdata} gives the QWK of EfficientNet-B3 and TOPORET on three datasets
acquired on different continents with different cameras~\cite{kaggle2015eyepacs,
decenciere2014messidor,porwal2018idrid}. The gain of TOPORET is larger out of
distribution ($+0.053$ and $+0.046$) than in distribution ($+0.040$), which is
consistent with TDA features being more invariant than CNN embeddings to acquisition
shifts. This would be an important property for screening, where images come from many
different camera models, but two external datasets are not enough to establish it in
general.

\begin{table}[H]
\centering
\caption[TOPORET cross-dataset results]{TOPORET cross-dataset results (OOD: out of distribution).}
\label{tab:crossdata}
\small
\begin{tabular}{@{}llrccc@{}}
\toprule
Dataset & Source & $N$ & EffNet-B3 & TOPORET & $\Delta$ \\
\midrule
EyePACS (in-dist.) & USA, multi-site & 88{,}702 & 0.840 & 0.880 & $+0.040$ \\
Messidor-2 (OOD) & France, other camera & 1{,}748 & 0.847 & 0.900 & $+0.053$ \\
IDRiD (OOD) & India, single site & 516 & 0.832 & 0.878 & $+0.046$ \\
\bottomrule
\end{tabular}
\end{table}

\subsection{Comparison with the Literature}
\label{sec:toporet-sota}

\Cref{tab:toporet-sota} compares TOPORET with the closest systems on the same benchmark.
Its QWK of $0.880$ is lower than that of RETFound ($0.920$); the difference is due to the
backbone (RETFound uses a ViT-L/16 pre-trained on $1.6$ million retinal images, TOPORET an
ImageNet EfficientNet-B3). This is by design: TOPORET was built to test H1 and H2 at
controlled backbone quality before adding RETFound in DOTS. Its contribution is not the
highest QWK but the joint treatment of L1 and L2 in a single framework.

\begin{table}[H]
\centering
\caption[TOPORET versus state-of-the-art systems (Kaggle EyePACS)]{TOPORET versus state-of-the-art systems (Kaggle EyePACS).}
\label{tab:toporet-sota}
\small
\begin{tabular}{@{}lccccc@{}}
\toprule
System & QWK & Sev-NR & L1 & L2 & Ordinal guarantee \\
\midrule
EfficientNet-B3 & 0.840 & 4.3\% & \xmark & \xmark & \xmark \\
RETFound & 0.920 & 3.1\% & \xmark & \xmark & \xmark \\
SAMF-GB & 0.882 & 3.8\% & \xmark & partial & \xmark \\
TOPORET & 0.880 & 3.4\% & \cmark & \cmark & \xmark \\
\bottomrule
\end{tabular}
\end{table}

\section{Discussion}
\label{sec:ch3-discussion}

\subsection{Why the Dual Metric Matters for H1}
\label{sec:dualnecessary}

A critical finding of \cref{tab:h1} is that the visual-only graph (GraphDR) fails the
pre-specified H1 threshold ($t_4=1.71$, $p=0.08$, n.s.), whereas the dual-metric graph
passes it ($t_4=3.82$, $p=0.009$). This is not a marginal difference: it separates a
non-significant gain from a very large confirmed effect. The Wasserstein topological
component is therefore, in these experiments, the component that makes the H1 test pass;
without it, the population graph provided no statistically significant benefit over
per-image classification.

\subsection{Biological Interpretation of the TDA Features}
\label{sec:tda-bio}

Each of the seven TDA features maps to a specific vascular pathology. $\Pi_{H_0}$ and
$n_{H_0}$ encode vascular rarefaction (they decrease with severity); $\Pi_{H_1}$,
$n_{H_1}$ and $\beta_1$ encode loop density from pathological anastomoses (they
increase); $D_f$ encodes the self-similarity of the vessel tree ($1.71\to1.57$ across
grades); and $\Delta$ (maximum $H_1$ persistence) captures neovascular fronds at Grade~4
($\eta^2=0.65$, the strongest effect). These mappings are post-hoc plausibility arguments
supporting H2; its statistical validity rests on the pre-specified ANOVA.

\subsection{From Supported Hypotheses to a Safety-Oriented System}
\label{sec:ch3-deploy}

H1 and H2 jointly support the view that the relational and structural axes contribute
useful information to grading. The Sev-NR of $3.4\%$ is a reminder that
statistical significance and clinical sufficiency are different standards. A $+4.0$~pp
QWK gain and a large ANOVA effect size are exactly what a methods paper would report as a
success; neither, alone or combined, certifies a system safe to triage real patients.

TOPORET's residual gap is therefore not a failure of H1 or H2 -- both remain
supported -- but a demonstration that confirming the relational and structural hypotheses
leaves the safety hypothesis (H3) untested. \Cref{tab:graphstruct} makes the mechanism
visible: even with the improved Grade-3 purity of $68.7\%$, roughly one in three of a
Severe-NPDR patient's nearest neighbours carries a different grade. This is exactly the
kind of residual ambiguity that an architecture without calibrated abstention cannot flag.
\Cref{ch:dots} is organised around closing this gap, together with two further gaps
(domain mismatch and ordinal rank violations) found by error analysis of the TOPORET
predictions.

\paragraph{Scope of H2: statistical association, not clinical validation.}
The ANOVA establishes that $H_0/H_1$ features are statistically associated with ICDR
grade with large effect sizes. Three caveats bound this claim.
\begin{enumerate}
\item \textbf{Association is not independent predictive value.} The ANOVA does not show
that TDA features add predictive value beyond visual features. The ablation of
\cref{tab:ablation3} provides this separately: TDA alone adds $+0.016$ QWK independently
of the visual backbone.
\item \textbf{Single-dataset validation.} H2 is tested on a stratified sample of Kaggle
EyePACS. Cross-dataset stability of the TDA feature distributions is supported by the
stability theorem (\cref{eq:stability}) but was not validated by an external ANOVA.
\item \textbf{Not a clinically validated biomarker.} Clinical validation would require
prospective studies, regulatory-grade reproducibility protocols and independent clinical
reading. The term used throughout this thesis is therefore \emph{severity-associated
topological feature}, not \emph{biomarker}.
\end{enumerate}

\subsection{Threats to Validity}
\label{sec:ch3-validity}

Three kinds of threats could weaken the conclusions of this chapter; each is stated here
together with the measure taken against it.

\paragraph{Internal validity.}
The main risk is that the gain attributed to the graph comes from another difference
between the systems, such as a different backbone or a different amount of tuning. All
three systems share the same frozen EfficientNet-B3 features, the same folds and the same
training budget, and the only difference between GraphDR and TOPORET is the similarity used
to build the graph. The hyperparameters were selected on training folds only, and the
decision thresholds of the tests were fixed before the experiments.

\paragraph{External validity.}
The hypotheses were tested on Kaggle EyePACS, a large but single benchmark whose labels come
from a screening network in the United States. The cross-dataset results on Messidor-2 and
IDRiD (\cref{tab:crossdata}) indicate that the gain persisted after a change of country and
camera, but the ANOVA of H2 was not repeated on external data, and the small size of IDRiD
limits the precision of the results obtained on it. All sample sizes are image counts;
some benchmarks contain several images of the same patient, so the results are image-level
and patient-level generalisation is interpreted cautiously.

\paragraph{Construct validity.}
QWK measures agreement with the reference grade, which itself contains grader noise
(\cref{sec:grading}); a part of the remaining disagreement may therefore reflect label
errors rather than model errors. The topological features depend on the quality of the
vessel segmentation, so a failure of stage S1 would appear as a change of topology.
Finally, H1 was tested with five folds, which is a small sample for a $t$-test; the
bootstrap interval and the exact permutation test reported with it were added for this reason;
with five folds, the permutation test cannot give a $p$-value below $1/32$.

\section*{Conclusion}
\phantomsection\addcontentsline{toc}{section}{Conclusion}

The pre-specified tests of TOPORET supported H1 ($t_4=3.82$, $p=0.009$, $d_z=1.71$, over
five folds) and H2 as a statistical association ($F(4,4995)=1308.9$, $p<0.001$,
$\eta^2=0.51$). The dual-similarity metric is a
prerequisite rather than an option: the visual-only graph gives a non-significant gain,
whereas the Wasserstein-augmented graph gives a very large one, with higher Grade-3
neighbourhood purity and larger out-of-distribution gains. The remaining gap -- Sev-NR
$=3.4\%$ against the $2\%$ threshold -- reflects the absence of a mechanism to defer
ambiguous cases, which the next chapter adds in the DOTS system.
